\documentclass[11pt,a4paper]{article}

% ── Packages ──────────────────────────────────────────────────────────────────
\usepackage[T1]{fontenc}
\usepackage[utf8]{inputenc}
\usepackage{palatino}
\usepackage{microtype}
\usepackage{geometry}
\usepackage{amsmath,amssymb,amsthm}
\usepackage{mathtools}
\usepackage{booktabs}
\usepackage{xcolor}
\usepackage[most]{tcolorbox}
\usepackage{hyperref}
\usepackage{fancyhdr}
\usepackage{titlesec}
\usepackage{parskip}
\usepackage{enumitem}

% ── Page layout ───────────────────────────────────────────────────────────────
\geometry{
  a4paper,
  top=2.75cm, bottom=3.15cm,
  left=2.75cm, right=2.75cm,
  headheight=22pt
}

% ── Colours ───────────────────────────────────────────────────────────────────
\definecolor{deepblue}{HTML}{1a1a2e}
\definecolor{midblue}{HTML}{2d3a7a}
\definecolor{lightblue}{HTML}{eef2ff}
\definecolor{ruleblue}{HTML}{9999bb}
\definecolor{boxframe}{HTML}{2d3a7a}
\definecolor{notegrey}{HTML}{555566}
\definecolor{accentgold}{HTML}{7a5c00}

% ── Hyperref ──────────────────────────────────────────────────────────────────
\hypersetup{
  colorlinks=true,
  linkcolor=midblue,
  urlcolor=midblue,
  citecolor=midblue,
  pdftitle={Causal Co-Scaling of Distance and Time in a Fractal S3 Reconstruction Model},
  pdfauthor={Andrei Preece, Boris Batenin},
  pdfsubject={Causal co-scaling, fractal S3 geometry, reconstruction time, speed of light}
}

% ── Section formatting ────────────────────────────────────────────────────────
\titleformat{\section}
  {\large\bfseries\color{deepblue}}
  {\thesection.}{0.6em}{}
  [\vspace{1pt}{\color{ruleblue}\hrule height 0.6pt}\vspace{2pt}]

\titleformat{\subsection}
  {\normalsize\bfseries\color{midblue}}
  {\thesubsection}{0.6em}{}

\titlespacing{\section}{0pt}{18pt}{8pt}
\titlespacing{\subsection}{0pt}{11pt}{4pt}

% ── Header / footer ───────────────────────────────────────────────────────────
\pagestyle{fancy}
\fancyhf{}
\fancyhead[L]{\small\color{midblue}\textit{Causal Co-Scaling — Fractal $S^3$ Geometry}}
\fancyhead[R]{\small\color{midblue}Preece \& Batenin, 2026}
\fancyfoot[C]{\small\color{notegrey}\thepage}
\renewcommand{\headrulewidth}{0.5pt}
\renewcommand{\headrule}{\hbox to\headwidth{%
  \color{ruleblue}\leaders\hrule height \headrulewidth\hfill}}

% ── Theorem styles ────────────────────────────────────────────────────────────
\newtheorem{axiom}{Axiom}
\newtheorem{proposition}{Proposition}
\newtheorem{remark}{Remark}

% ── tcolorbox styles ──────────────────────────────────────────────────────────
\newtcolorbox{principlebox}[1][]{
  enhanced, breakable,
  colback=deepblue!5,
  colframe=deepblue,
  boxrule=1pt, arc=4pt,
  left=10pt, right=10pt, top=8pt, bottom=8pt,
  title={\small\bfseries Causal Co-Scaling Principle},
  fonttitle=\bfseries\color{white},
  coltitle=white,
  attach boxed title to top left={yshift=-2.5mm, xshift=7mm},
  boxed title style={colback=deepblue, arc=3pt, boxrule=0pt, left=4pt, right=4pt},
  #1
}

\newtcolorbox{noteBox}[1][]{
  enhanced, breakable,
  colback=lightblue,
  colframe=boxframe,
  boxrule=0.75pt, arc=4pt,
  left=10pt, right=10pt, top=8pt, bottom=8pt,
  #1
}

\newtcolorbox{cautionBox}[1][]{
  enhanced, breakable,
  colback=accentgold!8,
  colframe=accentgold,
  boxrule=0.75pt, arc=4pt,
  left=10pt, right=10pt, top=8pt, bottom=8pt,
  #1
}

% ── Shortcuts ─────────────────────────────────────────────────────────────────
\newcommand{\Sthree}{S^{3}}
\newcommand{\Szero}{S^{0}}
\newcommand{\lP}{\ell_{P}}
\newcommand{\tP}{t_{P}}
\newcommand{\ph}{\varphi}
\newcommand{\Rrec}{R_{\mathrm{rec}}}
\newcommand{\Rstar}{R_{\ast}}
\newcommand{\tstar}{t_{\ast}}
\newcommand{\Vthree}{V_{S^3}}
\newcommand{\dd}{\mathrm{d}}

% ─────────────────────────────────────────────────────────────────────────────
\begin{document}
\thispagestyle{empty}

% ── Title block ───────────────────────────────────────────────────────────────
\begin{center}
  {\LARGE\bfseries\color{deepblue}
    Causal Co-Scaling of Distance and Time\\[5pt]
    in a Fractal $\Sthree$ Reconstruction Model}\\[8pt]
  {\Large\color{midblue}A Research Note on the Invariant Ratio $R/t=c$}\\[18pt]
  {\large Andrei Preece \quad \& \quad Boris Batenin}\\[7pt]
  {\small Independent Researchers}\\[4pt]
  {\small
    \href{mailto:andre.preece@web.de}{andre.preece@web.de}\quad
    ORCID: \href{https://orcid.org/0009-0004-5978-795X}{0009-0004-5978-795X}\\[2pt]
    \href{mailto:batenin@web.de}{batenin@web.de}\quad
    ORCID: \href{https://orcid.org/0009-0000-7027-8583}{0009-0000-7027-8583}\\[7pt]
    Berlin, 2026
  }
\end{center}

\vspace{4pt}
{\color{midblue}\hrule height 1.4pt}
\vspace{12pt}

\begin{center}
\begin{minipage}{0.93\linewidth}
\textbf{Manuscript scope.}\;
This research note isolates one structural axiom used in the broader compact $\Sthree$
programme: distance, reconstruction time, and the speed of light are not independent
quantities, but co-scaled components of one causal package.
\end{minipage}
\end{center}

\vspace{8pt}

% ── Abstract ──────────────────────────────────────────────────────────────────
\begin{center}
\begin{minipage}{0.93\linewidth}
\textbf{Abstract.}\;
This note formulates an axiom of \emph{causal co-scaling} for fractal compact
$\Sthree$ reconstruction.  The speed of light is not treated as an isolated number
to be derived after distance and time have already been defined.  Instead, the causal
ratio
\[
  c=\frac{R}{t}
\]
is taken as the invariant relation binding distance and reconstruction time.  If a
Planck-normalised hierarchy scales carrier radii as $R_N=\lP\ph^N$, then causal
self-consistency requires reconstruction times to scale as $t_N=\tP\ph^N$.  The ratio
is then preserved at every level:
\[
  \frac{R_N}{t_N}=\frac{\lP}{\tP}=c.
\]
This does not derive the numerical value of $c$, since both $\lP$ and $\tP$ already
encode the empirical causal scale.  Rather, it states a closure condition: a
golden-ratio reconstruction hierarchy is causally admissible only when distance and
time co-scale.  The note also separates the pre-photon reconstruction regime from
the later photon-supported cosmological expansion regime, and clarifies the status
of an internal transition index such as $N_{\ast}=207$.  Such an index is not an
observational definition of recombination, but a model-internal reference scale at
which the reconstruction mechanism is allowed to change regime.

\smallskip
\noindent\textbf{Keywords:} causal co-scaling;\ compact $\Sthree$;\;golden-ratio
hierarchy;\;Planck scale;\;reconstruction time;\;speed of light;\;fixed-point closure.
\end{minipage}
\end{center}

\vspace{10pt}
{\color{ruleblue}\hrule}
\vspace{14pt}

% ─────────────────────────────────────────────────────────────────────────────
\section{Purpose and Claim Status}
% ─────────────────────────────────────────────────────────────────────────────

The purpose of this note is narrow.  We isolate a single structural axiom:

\begin{quote}
\emph{If a carrier-distance scale is multiplied by a reconstruction factor, the
corresponding reconstruction time must be multiplied by the same factor, or else
the causal ratio $R/t$ is not preserved.}
\end{quote}

The note therefore does \textbf{not} claim:
\begin{enumerate}[leftmargin=2em, itemsep=3pt]
  \item to derive the empirical value of $c$ from pure topology;
  \item to derive the Planck length independently of $c$;
  \item to identify a fractal index with a directly observed cosmological epoch;
  \item to replace the standard observational definition of recombination.
\end{enumerate}

It claims only the following.

\begin{principlebox}
Distance, time, and the speed of light form a single causal triad:
\[
  R=ct,\qquad c=\frac{R}{t},\qquad t=\frac{R}{c}.
\]
In a reconstruction hierarchy, one cannot scale distance while leaving time unchanged.
Causal admissibility requires co-scaling.
\end{principlebox}

% ─────────────────────────────────────────────────────────────────────────────
\section{The Causal Triad}
% ─────────────────────────────────────────────────────────────────────────────

Let $R$ denote an admissible causal distance, $t$ the corresponding reconstruction
time, and $c$ the invariant causal speed.  The three elementary relations
\[
  R=ct,\qquad c=\frac{R}{t},\qquad t=\frac{R}{c}
\]
are formally equivalent, but conceptually they carry different emphases.

If $c$ is viewed as the primitive constant, then $R$ is the distance accessible during
time $t$.  If $R$ and $t$ are measured operationally, then $c$ appears as their invariant
ratio.  In a geometric reconstruction model, neither view alone is sufficient: the
causal geometry is the whole triad.

\begin{noteBox}
\textbf{Interpretation.}
The speed of light is not introduced here as an isolated object.  It is the causal
ratio binding distance and reconstruction time.  The model therefore asks not
``Can $c$ be derived from $\Sthree$ topology alone?'', but rather:
``Does the proposed $\Sthree$ reconstruction preserve the causal ratio $R/t=c$?''
\end{noteBox}

% ─────────────────────────────────────────────────────────────────────────────
\section{Planck Normalisation and the No-Derivation Caveat}
% ─────────────────────────────────────────────────────────────────────────────

The Planck length and Planck time are
\[
  \lP=\sqrt{\frac{\hbar G}{c^3}},
  \qquad
  \tP=\frac{\lP}{c}
  =\sqrt{\frac{\hbar G}{c^5}}.
\]
Thus any model that begins with $\lP$ or $\tP$ has already admitted the causal scale
$c$ into its normalisation.  Consequently, such a model cannot honestly derive $c$
as an independent result.

\begin{cautionBox}
\textbf{Caveat.}
Using $\lP$ and then obtaining $c=\lP/\tP$ is not a derivation of the speed of light.
It is a consistency check.  The value of $c$ is already present in the Planck
normalisation.
\end{cautionBox}

The correct target is therefore \emph{causal closure}.  A reconstruction hierarchy
must preserve the causal relation once the Planck causal unit has been fixed.

% ─────────────────────────────────────────────────────────────────────────────
\section{Axiom of Causal Co-Scaling}
% ─────────────────────────────────────────────────────────────────────────────

We now introduce the central axiom.

\begin{axiom}[Causal co-scaling]
Let a reconstruction hierarchy be indexed by $N\in\mathbb{Z}_{\ge 0}$ and let
$\ph=(1+\sqrt5)/2$.  If the carrier radius is scaled from the Planck scale by
\[
  R_N=\lP\ph^N,
\]
then the corresponding reconstruction time must scale as
\[
  t_N=\tP\ph^N.
\]
The causal ratio is then invariant:
\[
  \frac{R_N}{t_N}=c.
\]
\end{axiom}

\begin{proposition}[Causal closure]
Under causal co-scaling, every golden-ratio reconstruction level has the same
causal ratio:
\[
  \frac{R_N}{t_N}
  =
  \frac{\lP\ph^N}{\tP\ph^N}
  =
  \frac{\lP}{\tP}
  =
  c.
\]
\end{proposition}

\noindent
The proof is algebraic, but the interpretation is structural.  The factor $\ph^N$
is not permitted to multiply distance alone.  If it multiplies the carrier-distance,
it must multiply the reconstruction-time.  The scaling factor cancels from the ratio,
leaving the causal boundary unchanged.

% ─────────────────────────────────────────────────────────────────────────────
\section{Compact \texorpdfstring{$\Sthree$}{S3} Carrier Volume}
% ─────────────────────────────────────────────────────────────────────────────

If the spatial carrier is a round compact $\Sthree$ of radius $R$, its volume is
\[
  \Vthree(R)=2\pi^2R^3.
\]
This must be distinguished from the Euclidean ball volume
\[
  V_{\mathrm{ball}}(R)=\frac{4}{3}\pi R^3.
\]
The first is the volume of the three-sphere as a compact three-dimensional manifold.
The second is the volume enclosed by an ordinary two-sphere in Euclidean three-space.

For a Planck-normalised golden hierarchy,
\[
  R_N=\lP\ph^N,
\]
the compact $\Sthree$ carrier volume becomes
\[
  \Vthree(N)=2\pi^2\lP^3\ph^{3N}.
\]
If a fixed carrier content $E_0$ is distributed across this carrier, then the associated
energy-density scale is
\[
  \rho_E(N)=\frac{E_0}{\Vthree(N)}
  =
  \frac{E_0}{2\pi^2\lP^3}\ph^{-3N}.
\]
Thus a radius hierarchy scaling as $\ph^N$ induces a density hierarchy scaling as
$\ph^{-3N}$.

% ─────────────────────────────────────────────────────────────────────────────
\section{Distance, Time, Density, and Reconstruction Capacity}
% ─────────────────────────────────────────────────────────────────────────────

The co-scaling axiom gives a clean separation between causal scaling and density
scaling:
\[
  R_N\sim\ph^N,\qquad
  t_N\sim\ph^N,\qquad
  \rho_E(N)\sim\ph^{-3N}.
\]

The first two relations preserve the causal ratio.  The third expresses geometric
dilution of a fixed carrier content across an expanding compact volume.  This permits
one to introduce a distinct quantity: the number of reconstructible cells or the
reconstructible distinguishability capacity.  At resolution $\epsilon$, a generic
capacity law may be written
\[
  C_N(\epsilon)\sim \left(\frac{R_N}{\epsilon}\right)^{d_{\mathrm{eff}}},
\]
where $d_{\mathrm{eff}}$ is an effective reconstruction dimension.  At the Planck
resolution $\epsilon=\lP$,
\[
  C_N(\lP)\sim \left(\frac{R_N}{\lP}\right)^{d_{\mathrm{eff}}}
  =
  \ph^{d_{\mathrm{eff}}N}.
\]

This gives two opposite monotonic tendencies:
\[
  \rho_E(N)\sim\ph^{-3N},
  \qquad
  C_N(\lP)\sim\ph^{d_{\mathrm{eff}}N}.
\]

\begin{noteBox}
Energy density and reconstruction capacity are not the same quantity.  Energy
density measures concentration of carrier content.  Reconstruction capacity measures
how many distinguishable cells, modes, or states can be supported at a chosen
resolution.  Expansion may decrease the first while increasing the second.
\end{noteBox}

% ─────────────────────────────────────────────────────────────────────────────
\section{Two Regimes: Pre-Photon Reconstruction and Photon-Supported Expansion}
% ─────────────────────────────────────────────────────────────────────────────

The co-scaling law is most useful when the model distinguishes two regimes.

\subsection{Pre-photon reconstruction regime}

Before a free-photon description is introduced, one may model the carrier by pure
Planck-to-fractal reconstruction:
\[
  R_N=\lP\ph^N,
  \qquad
  t_N=\tP\ph^N.
\]
This is a structural hierarchy.  It is not yet an observational expansion law with a
measured Hubble parameter.  It states how causal distance and reconstruction time
must scale if the hierarchy is to preserve the invariant ratio $c$.

\subsection{Photon-supported expansion regime}

After the model introduces freely propagating photons, the description must change.
A later cosmological regime may involve quantities such as
\[
  a(t),\qquad H(t),\qquad z,\qquad R(t),
\]
and the empirical causal propagation of light.  The point of transition between the
two descriptions is not defined by the integer alone.  It requires a matching
condition.

% ─────────────────────────────────────────────────────────────────────────────
\section{Internal Transition Index}
% ─────────────────────────────────────────────────────────────────────────────

Suppose a model chooses an internal index $N_\ast$ and sets
\[
  \Rstar=\lP\ph^{N_\ast},
  \qquad
  \tstar=\tP\ph^{N_\ast}.
\]
Then by causal co-scaling,
\[
  \frac{\Rstar}{\tstar}=c.
\]
For example, if one takes $N_\ast=207$, this gives
\[
  \Rstar=\lP\ph^{207},
  \qquad
  \tstar=\tP\ph^{207}.
\]
This does not mean that ``the integer 207 is recombination.''  It means only that
$207$ is an internal reconstruction index, selecting a reference scale inside the
model.

\begin{cautionBox}
\textbf{Recombination caution.}
A physical recombination epoch is defined by physical criteria: photon decoupling,
temperature, ionisation state, redshift, and optical depth.  An internal golden-ratio
index may be used as a reference scale, but it does not define recombination by
itself.  To connect $N_\ast$ to recombination, one must supply an explicit matching
condition.
\end{cautionBox}

\begin{cautionBox}
\textbf{Numerical caution.}
Under the pure Planck normalisation of this note, $N_\ast=207$ evaluates to
\[
  \Rstar=\lP\ph^{207}\approx 2.94\times10^{8}\ \mathrm{m},
  \qquad
  \tstar=\tP\ph^{207}\approx 0.98\ \mathrm{s},
\]
i.e.\ a scale of about one light-second---\emph{not} a recombination scale.  A
light-travel time of order $3.8\times10^{5}$ years would instead require
$N\approx269.6$ under the same simple normalisation.  Hence $N_\ast=207$ can only be
an internal reference index of the reconstruction mechanism; any identification with
a cosmological epoch requires an explicit, separately stated matching condition, and
none is supplied or claimed here.
\end{cautionBox}

A precise statement is therefore:
\begin{quote}
If the pre-photon reconstruction mechanism is modeled by golden-ratio scaling from
the Planck carrier, then after $N_\ast$ reconstruction steps the radius reaches the
model's chosen pre-photon reference scale.  This does not observationally define
recombination; it marks a possible point at which the model changes regime from
$\ph$-reconstruction to photon-supported expansion.
\end{quote}

% ─────────────────────────────────────────────────────────────────────────────
\section{Relation to Previous Causal Self-Consistency Work}
% ─────────────────────────────────────────────────────────────────────────────

A previous draft on causal self-consistency treated the speed of light as an empirical
invariant and asked whether a fractal $\Sthree$ hierarchy could close consistently
under that invariant.  The present note extracts the simpler algebraic core of that
argument.

The earlier closure form may be summarised schematically as a fixed-point condition:
one accepts $c$ as an empirical causal speed and tests whether the proposed hierarchy
returns the same value after its internal geometric substitutions.  The present note
sharpens that idea:

\[
  R_N=\lP\ph^N,\qquad
  t_N=\tP\ph^N
  \quad\Longrightarrow\quad
  R_N/t_N=c.
\]

Thus the closure is no longer presented as a derivation of $c$, but as the
preservation of the causal triad under co-scaling.

% ─────────────────────────────────────────────────────────────────────────────
\section{Non-Claims}
% ─────────────────────────────────────────────────────────────────────────────

For clarity, the note explicitly rejects five overstatements.

\begin{enumerate}[leftmargin=2em, itemsep=3pt]
  \item \textbf{No derivation of $c$.}  Since $\lP$ and $\tP$ contain $c$, the ratio
  $\lP/\tP=c$ is a consistency identity, not an independent prediction.
  \item \textbf{No topology-only physics.}  Compact $\Sthree$ geometry supplies a
  carrier structure, not a complete derivation of empirical constants.
  \item \textbf{No direct recombination proof.}  A golden-ratio index requires a
  matching condition before it can be related to photon decoupling.
  \item \textbf{No replacement of observational cosmology.}  The framework is a
  reconstruction model; observational cosmology remains the test.
  \item \textbf{No observational support for the $\ph$-hierarchy is claimed.}  The
  directly testable log-periodic signature $\omega=2\pi/\ln\ph$ associated with
  golden-ratio scaling has been tested by us against DESI~DR2 BAO and Pantheon$+$
  data with null results,\footnote{B.~Batenin and A.~Preece, \emph{A Null Test of
  Golden-Ratio Log-Periodicity ($\omega=2\pi/\ln\varphi$) against DESI DR2 BAO and
  Pantheon$+$}, research manuscript (2026).}  and that status is inherited here: the
  co-scaling axiom is a structural admissibility condition on a model hierarchy, not
  an empirically supported scaling law.  Nothing in the DRD information-theory line
  depends on this note.
\end{enumerate}

% ─────────────────────────────────────────────────────────────────────────────
\section{Conclusion}
% ─────────────────────────────────────────────────────────────────────────────

The central result is the axiom of causal co-scaling:
\[
  R_N=\lP\ph^N,\qquad
  t_N=\tP\ph^N,\qquad
  \frac{R_N}{t_N}=c.
\]
It expresses the simple but structurally important fact that distance and
reconstruction time cannot be scaled independently once a causal speed is fixed.
The golden-ratio factor may organise a reconstruction hierarchy, but it does not
alter the causal ratio if time and distance co-scale.

In the broader compact $\Sthree$ programme, this axiom supplies a clean interface
between three ideas:
\[
  \text{Planck carrier}
  \quad\longrightarrow\quad
  \text{fractal reconstruction}
  \quad\longrightarrow\quad
  \text{photon-supported causal expansion}.
\]
The principle is modest, but useful.  It prevents false derivations of $c$ and replaces
them with a sharper demand: every proposed reconstruction scale must preserve the
causal triad of distance, time, and light.

\end{document}
